\input{@resources/tex-add-search-paths}\documentclass[SolvingMicroDSOPs]{subfiles}
\input{subfile-start}

\begin{document}
\hypertarget{the-infinite-horizon}{}
\section{The Infinite Horizon}\label{sec:the-infinite-horizon}

All of the solution methods presented so far have involved period-by-period iteration from an assumed last period of life, as is appropriate for life cycle problems.  However, if the parameter values for the problem satisfy certain conditions (detailed in \cite{BufferStockTheory}), the consumption rules (and the rest of the problem) will converge to a fixed rule as the horizon (remaining lifetime) gets large, as illustrated in Figure~\ref{fig:PlotCFuncsConverge}.  Furthermore, Deaton~\citeyearpar{deatonLiqConstr}, Carroll~\citeyearpar{carroll:brookings,carrollBSLCPIH} and others have argued that the `buffer-stock' saving behavior that emerges under some further restrictions on parameter values is a good approximation of the behavior of typical consumers over much of the lifetime.  Methods for finding the converged functions are therefore of interest, and are dealt with in this section.

Of course, the simplest such method is to solve the problem as
specified above for a large number of periods.  This is feasible, but
there are much faster methods.

\hypertarget{convergence}{}
\subsection{Convergence}\label{subsec:convergence}

In solving an infinite-horizon problem, it is necessary to have some
metric that determines when to stop because a solution that is `good
enough' has been found.

A natural metric is defined by the unique `target' level of wealth that \cite{BufferStockTheory} proves
will exist in problems of this kind \href{https://llorracc.github.io/BufferStockTheory\#GICNrm}{under certain conditions}: The $\mTrgNrm$ such that
\hypertarget{eq-mTrgNrmet}{}
\begin{equation}
  \Ex_{\prdt} [{\mNrm}_{\prdt+1}/\mNrm_{\prdt}] = 1 \text{~if~} \mNrm_{\prdt} = \mTrgNrm  \label{eq:mTrgNrmet}
\UnifiedNote{convergence criteria — steady-state target m̂ of the normalized MDP where 𝔼[m'/m] = 1; iteration stops when |m̂_{t+1} − m̂_t| < ε}
\end{equation}
where the accent is meant to signify that this is the value
that other $\mNrm$'s `point to.'

Given a consumption rule $\cFunc(\mNrm)$ it is straightforward to find
the corresponding $\mTrgNrm$.  So for our problem, a solution is declared
to have converged if the following criterion is met:
$\left|\mTrgNrm_{\prdt+1}-\mTrgNrm_{\prdt}\right| < \epsilon$, where $\epsilon$ is
a very small number and defines our degree of convergence tolerance.

Similar criteria can obviously be specified for other problems.
However, it is always wise to plot successive function differences and
to experiment a bit with convergence criteria to verify that the
function has converged for all practical purposes.

\begin{comment} % at suggestion of WW, this section was removed as unnecessary for the current model, which solves for the converged rule very fast
  \subsection{The Last Period}

  For the last period of a finite-horizon lifetime, in the absence of a
  bequest motive it is obvious that the optimal policy is to spend
  everything.  However, in an infinite-horizon problem there is no last
  period, and the policy of spending everything is obviously very far
  from optimal.  Generally speaking, it is much better to start off with
  a `last-period' consumption rule and value function equal to those
  corresponding to the infinite-horizon solution to the perfect
  foresight problem (assuming such a solution is known).

  For the perfect foresight infinite horizon consumption problem,
  the solution is
  \begin{equation}\begin{gathered}\begin{aligned}
        \bar{\cFunc}(m_{\prdt})  & = \overbrace{(1-\Rfree^{-1}(\Rfree
          \DiscFac)^{1/\CRRA})}^{\equiv
          \Min{\MPC}}\left[{m}_{\prdt}-1+\left(\frac{1}{1-1/\Rfree}\right)\right]
        \label{eq:pfinfhorc}
\UnifiedNote{[no direct counterpart] — perfect foresight consumption policy c̄(m) used as initial guess for the converged decision rule 𝒱(xᵥ)}
      \end{aligned}\end{gathered}\end{equation}
  where $\Min{\MPC}$ is the MPC in the
  infinite-horizon perfect foresight problem.  In our baseline problem,
  we set $\PermGroFac = \pLvl_{\prdt} = 1$.  It is straightforward to show that the
  infinite-horizon perfect-foresight value function and marginal value
  function are given by
  \begin{equation}\begin{gathered}\begin{aligned}
        \bar{\vFunc}(m_{\prdt})
        & =                                 \left(\frac{\bar{\cFunc}(m_{\prdt})^{1-\CRRA}}{
            (1-\CRRA)\Min{\MPC} }\right)
        \\  \bar{\vFunc}^{m}(m_{\prdt})  & =       (\bar{\cFunc}(m_{\prdt}))^{-\CRRA}
        \\  \Opt{\vFunc}^{m}(a_{\prdt})  & = \DiscFac \Rfree \PermGroFac_{\prdt+1}^{-\CRRA} \bar{\vFunc}^{m}(\RNrmByG_{\prdt+1} a_{\prdt}+1).
\UnifiedNote{[no direct counterpart] — perfect foresight 𝒱̄, 𝒱̄ᵐ, and continuation ℰ̄ᵐ(a) used to initialize value iteration}
      \end{aligned}\end{gathered}\end{equation}

\end{comment}
\end{document}
